\begin{corollary}[Standard-family nonvanishing for \texorpdfstring{$\operatorname{Pic}^0$}{Pic0}]
\hypertarget{tag-03DS}{}\label{wbII:cor:candidate-pic0-route}
\StatementRecord{wbII:cor:candidate-pic0-route}{proved}{03DS}{\SeriesRef{wbII:prop:pic0-pullback}, \SeriesRef{wbII:prop:conductor}, \SeriesRef{wbII:prop:cartier-duality-absolute}, \SeriesRef{wbII:lem:curve-flatness}, \SeriesRef{wbII:prop:analytic-base-change}}

For the standard threefold and every
$L\in\operatorname{Pic}^0(X)$,
\[
 H^0(W,\Omega_X^1|_W\otimes A_L)\ne0,
 \qquad
 H^2(W,(T_X\otimes L)|_W)\ne0,
\]
and
\[
 R^2f_*(T_X\otimes L)\ne0.
\tag{W.44}\hypertarget{tag-03GP}{}\label{wbII:eq:W:5.17f}
\]
\end{corollary}



\begin{proof}
The full cusp fibre $W=f^{-1}(p_0)$ is reduced by
\CompendiumRef{I}{def:common-conventions}.
At $k=l=1$, the cusp charts in the proof of
\CompendiumRef{I}{push:src:E.3.1} are $q=z_0z_1z_2$,
with the two-factor and smooth charts on the other strata,
and the normalization is $\eta:\mathrm{dP}_6\to W$.
The canonical-bundle isomorphism
$\omega_X\simeq\mathcal O_X(-2S_2)$ in
\CompendiumRef{I}{fam:src:A.1.1} gives a trivialization on
$X\setminus S_2$, which contains $W$: the two supports
are fibres over the distinct points $p_0$ and $p_2$. By
\cref{wbII:prop:pic0-pullback}, $\eta^*A_L$ is trivial, so
\cref{wbII:prop:conductor} gives the first nonvanishing.  Apply
\cref{wbII:prop:cartier-duality-absolute} to $E=T_X\otimes L$; its dual space is
exactly the preceding nonzero section space, giving the fibrewise $H^2$
nonvanishing.  The fibre has complex dimension two, so its coherent
$H^3$ vanishes.  The global construction makes $X$ compact and $f$
nonconstant; \cref{wbII:lem:curve-flatness} therefore gives properness and
$f$-flatness of $T_X\otimes L$.  The base-change isomorphism of
\cref{wbII:prop:analytic-base-change} identifies the fibre of the $R^2$
stalk with the nonzero $H^2$ group.  Hence the stalk, and therefore the
direct-image sheaf, is nonzero.

\end{proof}